% ============================================
% 数据流图模板 (Data Flow Diagram)
% 专业学术风格 - 适用于展示数据在系统中的流动
% ============================================
% Usage: Replace <caption>, <label>, and <nodes> placeholders
% Note: Colors (df-*) are defined in the main document preamble (survey.ini.tex)
%
% Key improvements:
% 1. Increased node widths to accommodate longer text
% 2. Support for multi-line text with proper wrapping
% 3. Better spacing between nodes
% 4. Improved flow label positioning
% 5. Dynamic height support

\begin{figure}[H]
\centering
\begin{adjustbox}{max width=0.95\textwidth, center}
\begin{tikzpicture}[
    node distance=2.0cm and 2.5cm,
    % 外部实体 - 红色方框 (用于输入/输出)
    external/.style={
        rectangle,
        draw=df-external!70,
        fill=df-external!15,
        minimum width=2.8cm,
        minimum height=1.4cm,
        text centered,
        font=\small\sffamily\bfseries,
        text=df-text,
        line width=1.2pt,
        text width=2.5cm,
        align=center,
        drop shadow={shadow xshift=0.5pt, shadow yshift=-0.5pt, opacity=0.1},
    },
    % 处理过程 - 蓝色圆角矩形
    process/.style={
        rectangle,
        rounded corners=8pt,
        draw=df-process!70,
        fill=df-process!15,
        minimum width=3.2cm,
        minimum height=1.4cm,
        text centered,
        font=\small\sffamily,
        text=df-text,
        line width=1pt,
        text width=3.0cm,
        align=center,
        drop shadow={shadow xshift=0.5pt, shadow yshift=-0.5pt, opacity=0.1},
    },
    % 宽处理过程 - 用于较长文本
    process-wide/.style={
        rectangle,
        rounded corners=8pt,
        draw=df-process!70,
        fill=df-process!15,
        minimum width=4.0cm,
        minimum height=1.4cm,
        text centered,
        font=\small\sffamily,
        text=df-text,
        line width=1pt,
        text width=3.8cm,
        align=center,
        drop shadow={shadow xshift=0.5pt, shadow yshift=-0.5pt, opacity=0.1},
    },
    % 数据存储 - 绿色开口矩形
    datastore/.style={
        rectangle,
        draw=df-datastore!70,
        fill=df-datastore!15,
        minimum width=3.2cm,
        minimum height=1.0cm,
        text centered,
        font=\small\sffamily,
        text=df-text,
        line width=1pt,
        text width=3.0cm,
        align=center,
    },
    % 宽数据存储
    datastore-wide/.style={
        rectangle,
        draw=df-datastore!70,
        fill=df-datastore!15,
        minimum width=4.0cm,
        minimum height=1.0cm,
        text centered,
        font=\small\sffamily,
        text=df-text,
        line width=1pt,
        text width=3.8cm,
        align=center,
    },
    % 数据流箭头
    dataflow/.style={
        ->,
        thick,
        >=stealth,
        draw=df-flow,
        line width=1.2pt,
    },
    % 数据流标签 - 改进版，支持多行
    flowlabel/.style={
        font=\footnotesize\sffamily,
        text=df-flow,
        fill=white,
        inner sep=3pt,
        align=center,
        text width=2.5cm,
    },
    % 宽流标签
    flowlabel-wide/.style={
        font=\footnotesize\sffamily,
        text=df-flow,
        fill=white,
        inner sep=3pt,
        align=center,
        text width=3.5cm,
    },
]

<nodes>

\end{tikzpicture}
\end{adjustbox}
\caption{<caption>}
\label{fig:<label>}
\end{figure}

% ============================================
% USAGE INSTRUCTIONS FOR CODE GENERATION:
% ============================================
%
% When generating dataflow diagrams, follow these guidelines:
%
% 1. NODE TYPES:
%    - external: For input/output entities (users, external systems)
%    - process: For processing steps (transformations, computations)
%    - process-wide: For processes with longer text
%    - datastore: For data storage (databases, files)
%    - datastore-wide: For datastores with longer text
%
% 2. NODE PLACEMENT (horizontal flow):
%    - For 6 nodes in a row: x = 0, 3.5, 7.0, 10.5, 14.0, 17.5
%    - For 5 nodes in a row: x = 0, 4.0, 8.0, 12.0, 16.0
%    - For 4 nodes in a row: x = 0, 4.5, 9.0, 13.5
%    - Minimum horizontal spacing: 3.5cm between node centers
%
% 3. TEXT HANDLING:
%    - Short text (< 12 chars): single line
%    - Medium text (12-20 chars): use text width for auto-wrap
%    - Long text (> 20 chars): use \\ for manual line breaks
%    - Example: {Query\\Augmentation} for two-line text
%
% 4. FLOW LABELS:
%    - Place labels at midpoint of arrows
%    - Use flowlabel for short labels
%    - Use flowlabel-wide for longer labels
%    - Use \\ for multi-line labels
%
% EXAMPLE - Linear data flow with 5 nodes:
%
% % Nodes
% \node[external] (input) at (0, 0) {Raw Query};
% \node[process] (proc1) at (4, 0) {Query\\Augmentation};
% \node[datastore-wide] (store) at (8, 0) {Tool Repository\\(API Metadata)};
% \node[process-wide] (proc2) at (12, 0) {Dense Vector\\Retrieval};
% \node[external] (output) at (16, 0) {Selected\\Candidates};
%
% % Flows with labels
% \draw[dataflow] (input) -- (proc1) node[flowlabel, midway, above] {raw query};
% \draw[dataflow] (proc1) -- (store) node[flowlabel, midway, above] {augmented\\vector};
% \draw[dataflow] (store) -- (proc2) node[flowlabel, midway, above] {technical\\docs};
% \draw[dataflow] (proc2) -- (output) node[flowlabel, midway, above] {final\\selection};
%
% EXAMPLE - Vertical data flow:
%
% \node[external] (input) at (0, 0) {User Input};
% \node[process] (proc1) at (0, -2.5) {Processing};
% \node[datastore] (store) at (0, -5) {Database};
%
% \draw[dataflow] (input) -- (proc1);
% \draw[dataflow] (proc1) -- (store);